\subsection{A conditional estimate using perpendicular bisectors}
\label{bisector:section}

The following criterion retains the positions of source pairs, which are
absent from the scale profile. Its mixed-norm hypothesis is not deduced
from Hausdorff and packing dimensions here. In particular it is not a new
dimensional sufficient condition.

Let \(\mu,\nu\) be compactly supported positive measures with
\(I_1(\mu)<\infty\) and \(I_\tau(\nu)<\infty\), where
\(1<\tau<2\). For \(x\ne x'\), define
\[
 \theta(x,x')=\frac{x-x'}{|x-x'|},\qquad
 c(x,x')=\frac{|x|^2-|x'|^2}{2|x-x'|}.
\]
Their perpendicular bisector has equation \(y\cdot\theta=c\).
The finite measure on oriented line space
\[
 \mathcal B_\mu=(\theta,c)_*
       \bigl(|x-x'|^{-1}\mu(dx)\mu(dx')\bigr)
\]
has total mass \(I_1(\mu)\). We use ordinary angular arclength below.

\begin{proposition}[Mixed-norm bisector criterion]\label{bisector:criterion}
Suppose
\[
 \mathcal B_\mu=b_\mu(\theta,c)\,d\theta\,dc,\qquad
 \|b_\mu\|_{L^2_\theta L^{2/\tau}_c}<\infty.
\]
Then the actual joint distance law has an \(L^2(\nu\times dt)\)
density. Writing \(\eta_y=(d_y)_*\mu\), one has
\begin{equation}\label{bisector:bound}
 \int\|\eta_y\|_2^2\,d\nu(y)
 \le C_{\tau,K}
   \bigl(\nu(\R^2)^2+I_\tau(\nu)\bigr)^{1/2}
   \|b_\mu\|_{L^2_\theta L^{2/\tau}_c},
\end{equation}
where \(K\) is a compact set containing both supports. Positive source
mass therefore gives positive distance length for almost every pin.
\end{proposition}

\begin{proof}
For almost every \(\theta\), let \(v_\theta\) be the density of the
linear projection of \(\nu\) in direction \(\theta\). The Fourier
projection identity and polar coordinates imply
\[
 \int\|v_\theta\|_{H^{(\tau-1)/2}(\R)}^2\,d\theta
 \le C_\tau\bigl(\nu(\R^2)^2+I_\tau(\nu)\bigr).
\]
Indeed the high-frequency weight is \(|\xi|^{\tau-2}\) in the
plane; the bounded-frequency part is controlled by the mass.
One-dimensional Sobolev embedding and the Hardy--Littlewood maximal
inequality give, for \(q=2/(2-\tau)\),
\begin{equation}\label{bisector:projection}
 \|Mv_\theta\|_{L^2_\theta L^q_c}
 \le C_\tau\bigl(\nu(\R^2)^2+I_\tau(\nu)\bigr)^{1/2}.
\end{equation}
Here \(M\) is the centered interval maximal operator and
\(q'=2/\tau\).

Choose a nonnegative smooth compactly supported line mollifier
\(\varphi_\varepsilon\) of integral one. Expanding its squared
convolution norm bounds
\(\int\|\eta_y*\varphi_\varepsilon\|_2^2\,d\nu(y)\) by
\[
 C\varepsilon^{-1}\iint
 \nu\{y:\bigl||x-y|-|x'-y|\bigr|\le C\varepsilon\}
 \,d\mu(x)d\mu(x').
\]
All distances are bounded above on \(K\), and
\[
 |x-y|^2-|x'-y|^2
 =2|x-x'|\bigl(c(x,x')-\theta(x,x')\cdot y\bigr).
\]
Thus the integrand is at most
\(C_K|x-x'|^{-1}Mv_\theta(c)\).
The exceptional angular null set has zero weighted source-pair mass
by the assumed density of \(\mathcal B_\mu\). The source diagonal
also has zero mass, because finite first energy excludes atoms.
Consequently the smoothed squared norm is bounded by
\[
 C_K\int b_\mu(\theta,c)Mv_\theta(c)\,dc\,d\theta.
\]
Apply H\"older with \(q',q\) in \(c\), followed by Cauchy--Schwarz
in \(\theta\) and \eqref{bisector:projection}. This proves
\eqref{bisector:bound} uniformly in \(\varepsilon\). Weak compactness
in the joint \(L^2\) space identifies an actual density by testing
against compactly supported continuous functions. Disintegration gives
the pinned statement. No separation of the two supports is needed.
\end{proof}

For a smooth source density \(f\), the bisector density has an exact
reflection formula:
\begin{equation}\label{bisector:reflection}
 b_\mu(\theta,c)=\int_{\R^2}f(x)f(\mathcal R_{\theta,c}x)\,dx,
 \qquad \mathcal R_{\theta,c}x=x+2(c-x\cdot\theta)\theta.
\end{equation}
To check the normalization, put
\(x=z+r\theta/2\), \(x'=z-r\theta/2\), and
\(z=c\theta+u\theta^\perp\), with \(r>0\). The pair Jacobian
\(r\) cancels the reciprocal separation. Changing \(r\) to twice
the normal displacement and using evenness of the product converts
the integral over positive displacements into the integral over the
whole normal line. This gives \eqref{bisector:reflection} with no
additional factor.

For an \(s\)-Frostman source, \(s>1\), a sufficient finite-resolution
version is the uniform bound
\[
 \sup_{0<\delta<1}
 \left\|\int f_\delta(x)f_\delta(\mathcal R_{\theta,c}x)\,dx
       \right\|_{L^2_\theta L^{2/\tau}_c}<\infty,
 \qquad f_\delta=\mu*\psi_\delta,
\]
with a fixed positive smooth compactly supported planar mollifier.
The smoothed measures retain the Frostman bound, also at radii below
\(\delta\). Their weighted near-diagonal mass is uniformly
\(O(r^{s-1})\). Away from the diagonal, the bisector map is continuous,
so these measures converge weakly to \(\mathcal B_\mu\).
Reflexivity of the mixed-norm space proves the stated implication.

For example, at \(d=23/20\), \(D=3/2\), which lies outside the
current cutoff \(B_{\rm H}(23/20)=10/7\), choose a pin energy exponent
\(\tau=9/8\). The source reflection norm would be
\(L^2_\theta L^{16/9}_c\), paired with
\(L^2_\theta L^{16/7}_c\) for pin projections. Proving the required
reflection bound for a positive source restriction under these
dimensions remains an open step in this approach. The criterion has
a raw quadratic conclusion, so the earlier counterexamples forbid
asserting its hypothesis for arbitrary Frostman measures.
